\documentclass[12pt, a4paper]{article}
\usepackage[utf8]{inputenc}
\usepackage[T1]{fontenc}
\usepackage[english]{babel}
\usepackage{geometry}
\geometry{top=2.5cm, bottom=2.5cm, left=2cm, right=2cm}
\usepackage{parskip}
\usepackage{microtype}
\usepackage[default]{sourcesanspro}
\usepackage{xcolor}
\usepackage{fancyhdr}
\usepackage[most]{tcolorbox}
\usepackage[hidelinks]{hyperref}
\usepackage{amssymb, amsmath}

\newcommand{\Book}{Book Title}
\newcommand{\Chapter}{Chapter Title}
\definecolor{StructureColor}{RGB}{40, 80, 120}

\pagestyle{fancy}
\fancyhf{}
\lhead{\color{StructureColor}\textbf{\Book}}
\rhead{\small\Chapter}
\cfoot{\thepage}
\renewcommand{\headrulewidth}{0.4pt}
\renewcommand{\headrule}{\hbox to\headwidth{\color{StructureColor}\leaders\hrule height \headrulewidth\hfill}}

\newtcolorbox{exercisebox}[1]{
    enhanced, colback=StructureColor!5!white, colframe=StructureColor,
    coltitle=white, fonttitle=\bfseries\sffamily,
    title={#1}, sharp corners, boxrule=0.5mm
}

\begin{document}
\thispagestyle{fancy}

\begin{exercisebox}{Probability Exercise}
    Let $X$ be a random variable such that $X \sim Bin(n, p)$. Find the mode of $X$. Justify your answer.
\end{exercisebox}

\textbf{Solution}

The mode of $X$ is the $k$ where

\[
P(X=k) \geq P(X=k+1) \land  P(X=k) \geq P(X=k-1)
\]

Start solving for $P(X=k) \geq P(X=k+1)$:
$$
\begin{aligned}
    P(X=k) &\geq P(X=k+1)\\\\
    \dfrac{P(X=k)}{P(X=k+1)}
    &\geq1\\\\
    \dfrac{\binom{n}{k}\cdot p^k \cdot (1-p)^{n-k}}
    {\binom{n}{k+1}\cdot p^{k+1} \cdot (1-p)^{n-(k+1)}}
    &\geq1\\\\
    \dfrac{\dfrac{n!}{k!\cdot (n-k)!}}
    {\dfrac{n!}{(k+1)!\cdot (n-(k+1))!}}
    \cdot 
    \dfrac{1-p}{p}
    &\geq1\\\\
    \dfrac{(k+1)}{(n-k)}
    \cdot 
    \dfrac{1-p}{p}
    &\geq1\\\\
    (k+1)(1-p) &\geq (n-k)\cdot p\\\\
    k - kp + 1 - p &\geq np - kp\\\\
    k &\geq p(n+1) - 1
\end{aligned}
$$

Now solving for $P(X=k) \geq P(X=k-1)$:

$$
\begin{aligned}
    P(X=k) &\geq P(X=k-1)\\\\
    \dfrac{P(X=k)}{P(X=k-1)}
    &\geq1\\\\
    \dfrac{\binom{n}{k}\cdot p^k \cdot (1-p)^{n-k}}
    {\binom{n}{k-1}\cdot p^{k-1} \cdot (1-p)^{n-(k-1)}}
    &\geq1\\\\
    \dfrac{\dfrac{n!}{k!\cdot (n-k)!}}
    {\dfrac{n!}{(k-1)!\cdot (n-k+1)!}}
    \cdot 
    \dfrac{p}{1-p}
    &\geq1\\\\
    \dfrac{(n-k+1)}{k}
    \cdot 
    \dfrac{p}{1-p}
    &\geq1\\\\
    (n-k+1)\cdot p &\geq k\cdot(1-p)\\\\
    np - kp + p &\geq k - kp\\\\
    k &\leq p(n+1)
\end{aligned}
$$



Therefore, the mode of $X$ is where $k$ is in the interval:

$$\boxed{p(n+1) - 1 \leq k \leq p(n+1)}$$

\end{document}